\documentclass[11pt]{article}

\usepackage{amsmath,amssymb,amsthm}
\usepackage{fullpage}
\usepackage{hyperref}

\title{Scoped Universality, Cycle–Overlap Rigidity, and the EntropyDepth Wall}
\author{Inacio F. Vasquez \\ Independent Researcher (Foundations of Physics \& Mathematics)}
\date{\today}

\newtheorem{definition}{Definition}
\newtheorem{lemma}{Lemma}
\newtheorem{theorem}{Theorem}
\newtheorem{corollary}{Corollary}

\begin{document}

\maketitle

\begin{abstract}
We formalize a corrected universality principle governing refinement-based SAT algorithms.
The key invariant is the cycle–overlap rank $\CR(G_F)$ of the incidence graph of a CNF formula.
We prove that bounded-degree CNF families with linear cycle–overlap rank force linear
EntropyDepth for every FO$^k$-admissible refinement process.  
This establishes a precise equivalence between three conditions: cycle richness,
FO$^k$-type diversity, and linear refinement depth.
\end{abstract}

\section{Incidence Graphs}

\begin{definition}[Incidence Graph]
Let $F$ be a CNF formula with variable set $X$ and clause set $\mathcal{C}$.

The \emph{incidence graph} is the bipartite graph

\[
G_F = (X \cup \mathcal{C},E)
\]

with edge $(x,C)$ whenever variable $x$ appears in clause $C$.
\end{definition}

\begin{definition}[Bounded Degree]
A CNF family $\{F_n\}$ has bounded degree $\Delta$ if

\[
\deg(G_{F_n}) \le \Delta
\]

for all $n$.
\end{definition}

---

\section{Cycle–Overlap Rank}

Let $Z_1(G,\mathbb{F}_2)$ denote the cycle space of $G$.

\begin{definition}[Cycle–Overlap Rank]
For a graph $G$, define

\[
\CR(G) =
\dim_{\mathbb{F}_2}
\mathrm{span}\{C_1\cap C_2 :
C_1,C_2 \in Z_1(G,\mathbb{F}_2)\}
\]

where intersections are taken as edge sets.
\end{definition}

Intuitively this measures the number of independent cycle interactions.

---

\section{CR–Rich Formula Families}

\begin{definition}[CR–Rich]
A CNF family $\{F_n\}$ is \emph{CR–Rich} if

\[
\CR(G_{F_n}) = \Omega(n)
\]
\end{definition}

This property holds for Tseitin formulas on expanders.

---

\section{FO$^k$ Local Types}

Fix $k$ variables and radius

\[
R = 3^k
\]

by Gaifman locality.

\begin{definition}[FO$^k$ Type]
Two vertices have the same FO$^k_R$ type if they satisfy the same
first-order formulas using $k$ variables within radius $R$.
\end{definition}

Let

\[
|\mathrm{FO}^k_R(G)|
\]

denote the number of such types.

---

\section{Cycle Rigidity}

\begin{lemma}[Cycle Rigidity]
Let $G$ be bounded-degree.

If

\[
\CR(G) = \Omega(n)
\]

then

\[
|\mathrm{FO}^k_R(G)| \ge 2
\]

for sufficiently large $n$.
\end{lemma}

\begin{proof}
If all vertices shared one FO$^k$ type, every radius-$R$
neighborhood would be isomorphic.

This would imply bounded cycle interaction inside each
radius-$R$ neighborhood.

Thus the global cycle–overlap rank would be bounded.

Contradiction.
\end{proof}

---

\section{EntropyDepth}

\begin{definition}[Refinement Process]
A refinement process $P$ iteratively splits variable sets using
local rules depending on bounded neighborhoods.
\end{definition}

\begin{definition}[EntropyDepth]
Let $P_n$ be the process on formula $F_n$.

\[
\mathrm{ED}(P_n)
\]

is the number of refinement steps required until all variables
are separated.
\end{definition}

---

\section{Linear EntropyDepth}

\begin{lemma}[Entropy Bound]
If

\[
|\mathrm{FO}^k_R(G)| \ge 2
\]

then every FO$^k$-admissible refinement process requires

\[
\mathrm{ED}(P_n) = \Omega(n).
\]
\end{lemma}

\begin{proof}
Each refinement step separates at most $O(1)$ new FO$^k$ types.

Thus separating $n$ variables requires $\Omega(n)$ steps.
\end{proof}

---

\section{Scoped Universality}

\begin{theorem}[Scoped Universality]
Let $\{F_n\}$ be bounded-degree and CR–Rich.

Then for every FO$^k$-admissible refinement process $P$

\[
\mathrm{ED}(P_n) = \Theta(n).
\]
\end{theorem}

\begin{proof}
Lower bound from the previous lemma.

Upper bound follows because each step separates
at least one variable.
\end{proof}

---

\section{Definitive Equivalence}

\begin{theorem}[Definitive Universality]
For bounded-degree CNF families the following are equivalent.

\begin{align}
\mathrm{(A)} & \quad \CR(G_{F_n}) = \Omega(n) \\
\mathrm{(B)} & \quad \text{CR–Rich property holds} \\
\mathrm{(C)} & \quad \forall P: \mathrm{ED}(P_n) = \Theta(n)
\end{align}

where $P$ ranges over FO$^k$-admissible refinement processes.
\end{theorem}

\begin{proof}
(A)$\Rightarrow$(B) immediate.

(B)$\Rightarrow$(C) from cycle rigidity and the entropy bound.

(C)$\Rightarrow$(A) because if $\CR(G_{F_n})$ were sublinear,
cycle interactions would remain bounded,
yielding bounded FO$^k$ types and sublinear depth.
\end{proof}

---

\section{Consequences}

For Tseitin formulas on expander graphs

\[
\CR(G) = \Theta(n)
\]

thus

\[
\mathrm{ED}(P_n) = \Theta(n)
\]

for every FO$^k$ refinement algorithm.

This establishes the EntropyDepth wall.

---

\section{Conclusion}

Cycle–overlap rank provides the missing invariant linking
graph topology, logical distinguishability, and refinement depth.

Linear cycle interaction forces linear algorithmic depth
for all locality-bounded SAT reasoning.

\end{document}

